\documentclass[en,hazy,blue,12pt,normal]{elegantnote}

\input{../../preface}

\title{DSC 206: Homework 2}
\author{Zeyu Bian}
\date{\today}

\begin{document}

\maketitle

%================================================================
\begin{homework}
 Let $P=\{p_1,\dots,p_n\}\subset\MR$ with $p_1\le p_2\le\cdots\le p_n$.
\begin{enumerate}
    \item[(1.a)] Describe the optimal $1$-center.
    \item[(1.b)] Describe the optimal $1$-mean.
    \item[(1.c)] Give an example where the $1$-median minimizer is not unique. Then, assuming $n$ is odd, find $p^*=\arg\min_{p\in P}\sum_{i=1}^n |p_i-p|$ and prove your answer.
\end{enumerate}
\end{homework}

\begin{proof}[\solutionname]

\textbf{(1.a) Optimal $1$-center.}
The $1$-center cost is $f(x)=\max_i |p_i-x|=\max(x-p_1,\,p_n-x)$ since the points are sorted. This is minimized when both terms are equal, i.e.,
\[
x^*=\frac{p_1+p_n}{2},\qquad \text{cost}=\frac{p_n-p_1}{2}.
\]

\textbf{(1.b) Optimal $1$-mean.}
The $1$-mean cost is $g(x)=\sum_{i=1}^n (p_i-x)^2$. Setting $g'(x)=-2\sum_{i=1}^n(p_i-x)=0$ yields
\[
x^*=\bar p=\frac{1}{n}\sum_{i=1}^n p_i.
\]
Since $g''(x)=2n>0$, this is the unique minimizer.

\textbf{(1.c) $1$-median.}

\emph{Non-uniqueness example.} Take $n=2$ with $P=\{0,1\}$. Then $h(x)=|x|+|x-1|=1$ for every $x\in[0,1]$, so every point in $[0,1]$ is a minimizer.

\emph{Odd $n$.} We will show $p^*=p_{(n+1)/2}$ is the median element of $P$.

\emph{Proof.} The function $h(x)=\sum_{i=1}^n|p_i-x|$ is convex and piecewise linear. On the open interval $(p_k,p_{k+1})$ for $1\le k\le n-1$ we have
\[
h'(x)=\#\{i:p_i<x\}-\#\{i:p_i>x\}=k-(n-k)=2k-n.
\]
For odd $n$, it is negative for $k\le \frac{n-1}{2}$ and positive for $k\ge \frac{n+1}{2}$. Thus $h$ strictly decreases on $(-\infty,p_{(n+1)/2}]$ and strictly increases on $[p_{(n+1)/2},\infty)$. Hence the unique minimizer over all of $\MR$ is
\[
p^*=p_{(n+1)/2}.\qedhere
\]
\end{proof}

%================================================================
\begin{homework}
Let $\mu^*$ be the optimal $k$-center cost (with arbitrary centers in $\MR^d$) and $\widehat\mu^*$ the optimal restricted $k$-center cost (centers must lie in $P$). Prove $\mu^*\le \widehat\mu^*\le 2\mu^*$.
\end{homework}

\begin{proof}[\solutionname]

\textbf{Lower bound $\mu^*\le \widehat\mu^*$.}
Every restricted clustering is also a feasible (unrestricted) $k$-clustering. Therefore the optimum over a smaller (restricted) feasible set is at least the optimum over the larger feasible set:
$\mu^*=\min_{\text{all}}\text{cost}\le \min_{\text{restricted}}\text{cost}=\widehat\mu^*.$

\textbf{Upper bound $\widehat\mu^*\le 2\mu^*$.}
Let $\{(c_i^*,C_i^*)\}_{i\in[1,k]}$ be an optimal unrestricted clustering with cost $\mu^*$. For each $i$, choose any data point $\widehat c_i\in C_i^*$ as a restricted center. Keep the partition be the same. For any $p\in C_i^*$, the triangle inequality gives
\[
d(p,\widehat c_i)\le d(p,c_i^*)+d(c_i^*,\widehat c_i)\le \mu^*+\mu^*=2\mu^*,
\]
where both bounds use the definition $d(q,c_i^*)\le \mu^*$ for all $q\in C_i^*$. Thus this restricted clustering has cost at most $2\mu^*$, and since $\widehat\mu^*$ is the minimum over all restricted clusterings,
\[
\widehat\mu^*\le 2\mu^*.\qedhere
\]
\end{proof}

%================================================================
\begin{homework}
Graph $G=(V,E)$ has three connected components on vertex sets $V_1=\{v_1,\dots,v_{n/3}\}$, $V_2=\{v_{n/3+1},\dots,v_{2n/3}\}$, $V_3=\{v_{2n/3+1},\dots,v_n\}$. Provide three orthonormal eigenvectors of the Laplacian $L$ for eigenvalue $0$.
\end{homework}

\begin{proof}[\solutionname]
For any connected component, the indicator vector of that component is a null vector of $L$, since $L\bm 1_V=0$ on each component (the row sums of $L$ restricted to a component are zero). The three indicator vectors of $V_1,V_2,V_3$ have disjoint support, so they are mutually orthogonal. Normalizing each to unit length (each has $n/3$ ones) gives the orthonormal basis
\[
\phi_i \;=\; \sqrt{\tfrac{3}{n}}\;\bm 1_{V_i},\qquad i=1,2,3,
\]
explicitly $(\phi_i)_j=\sqrt{3/n}$ if $v_j\in V_i$ and $0$ otherwise. By construction $L\phi_i=0$, $\|\phi_i\|=1$, and $\langle\phi_i,\phi_j\rangle=0$ for $i\ne j$.
\end{proof}

%================================================================
\begin{homework}
Let $S$ be a finite point set with centroid $\mu(S)$, and let $T$ be disjoint from $S$. Show that
\[
\bigl\|\mu(S\cup T)-\mu(S)\bigr\|\le \frac{|T|}{|S|+|T|}\,\bigl\|\mu(T)-\mu(S)\bigr\|.
\]
\end{homework}

\begin{proof}[\solutionname]
Write $|S|=m$ and $|T|=t$. By definition,
\[
\mu(S\cup T)=\frac{1}{m+t}\sum_{x\in S\cup T} x=\frac{1}{m+t}\Bigl(\sum_{x\in S}x+\sum_{x\in T}x\Bigr)=\frac{m\,\mu(S)+t\,\mu(T)}{m+t}.
\]
Therefore
\[
\mu(S\cup T)-\mu(S)=\frac{m\,\mu(S)+t\,\mu(T)}{m+t}-\mu(S)=\frac{t\bigl(\mu(T)-\mu(S)\bigr)}{m+t},
\]
and taking norms,
\[
\bigl\|\mu(S\cup T)-\mu(S)\bigr\|=\frac{t}{m+t}\bigl\|\mu(T)-\mu(S)\bigr\|=\frac{|T|}{|S|+|T|}\bigl\|\mu(T)-\mu(S)\bigr\|.
\]
The inequality follows.
\end{proof}

%================================================================
\begin{homework}
For the weighted graph in the assignment:
\begin{enumerate}
    \item[(5.a)] Construct the graph Laplacian $L$.
    \item[(5.b)] Compute the bipartition obtained by spectral clustering: give the second eigenvector and the induced partition.
\end{enumerate}
\end{homework}

\begin{proof}[\solutionname]

\textbf{(5.a) Laplacian.} Reading the edge weights from the figure: heavy edges $(v_1,v_2)=50$, $(v_1,v_3)=50$, $(v_4,v_5)=50$; light edges $(v_1,v_4)=1$, $(v_1,v_5)=1$, $(v_2,v_4)=1$, $(v_2,v_5)=1$, $(v_3,v_5)=1$. The weighted degrees are
\[
d_1=102,\;d_2=52,\;d_3=51,\;d_4=52,\;d_5=53.
\]
Thus $L=D-W$ equals
\[
L=
\begin{pmatrix}
102 & -50 & -50 & -1 & -1\\
-50 & 52 & 0 & -1 & -1\\
-50 & 0 & 51 & 0 & -1\\
-1 & -1 & 0 & 52 & -50\\
-1 & -1 & -1 & -50 & 53
\end{pmatrix}.
\]

\textbf{(5.b) Spectral bipartition.}
The smallest eigenvalue of $L$ is $0$ with eigenvector $\bm 1/\sqrt{5}$. The second-smallest eigenvalue vector $\phi_2$ approximately separates the two clusters of strongly-connected vertices. Heavy edges of weight $50$ form two tightly bound groups: $A=\{v_1,v_2,v_3\}$ joined through $v_1$, and $B=\{v_4,v_5\}$ joined directly. Light weight-$1$ edges only weakly connect $A$ and $B$.

The Fiedler vector takes opposite signs on $A$ and $B$. To leading order it is the indicator-difference vector orthogonal to $\bm 1$. With $|A|=3$, $|B|=2$, requiring $\sum_i\phi_2(i)=0$ and $\|\phi_2\|=1$ gives
\[
\phi_2 \;\approx\; \Bigl(\sqrt{\tfrac{2}{15}},\;\sqrt{\tfrac{2}{15}},\;\sqrt{\tfrac{2}{15}},\;-\sqrt{\tfrac{3}{10}},\;-\sqrt{\tfrac{3}{10}}\Bigr)
.
\]
The induced bipartition by sign of $\phi_2$ is
\[
C_1=\{v_1,v_2,v_3\},\qquad C_2=\{v_4,v_5\}.
\]
This matches intuition: the spectral cut severs only weight-$1$ edges and preserves all weight-$50$ edges.
\end{proof}

%================================================================
\begin{homework}
\textbf{Problem 6 (3 pts).} Compute the Jaccard similarities of each pair of $\{1,2,3,4\}$, $\{2,3,5,7\}$, $\{2,4,6\}$.
\end{homework}

\begin{proof}[\solutionname]
Let $A=\{1,2,3,4\},\,B=\{2,3,5,7\},\,C=\{2,4,6\}$. Recall $J(X,Y)=|X\cap Y|/|X\cup Y|$.
\begin{align*}
J(A,B)&=\frac{|\{2,3\}|}{|\{1,2,3,4,5,7\}|}=\frac{2}{6}=\frac{1}{3},\\[2pt]
J(A,C)&=\frac{|\{2,4\}|}{|\{1,2,3,4,6\}|}=\frac{2}{5},\\[2pt]
J(B,C)&=\frac{|\{2\}|}{|\{2,3,4,5,6,7\}|}=\frac{1}{6}.
\end{align*}
\end{proof}

%================================================================
\begin{homework}
Using HW figure, augment the signatures with $h_3(x)=2x+1\bmod 5$ and $h_4(x)=3x-1\bmod 5$.
\end{homework}

\begin{proof}[\solutionname]

The hash values per row are
\[
\begin{array}{c|ccccc}
x & 0 & 1 & 2 & 3 & 4\\\hline
h_3(x)=2x+1\bmod 5 & 1 & 3 & 0 & 2 & 4\\
h_4(x)=3x-1\bmod 5 & 4 & 2 & 0 & 3 & 1
\end{array}
\]
For each column $S_j$, the minhash with hash $h$ is $\min\{h(r):S_j(r)=1\}$.

\begin{itemize}
    \item $S_1$ has $1$'s at rows $\{0,3\}$: $h_3=\min(1,2)=1$, $h_4=\min(4,3)=3$.
    \item $S_2$ has $1$'s at rows $\{2\}$: $h_3=0$, $h_4=0$.
    \item $S_3$ has $1$'s at rows $\{1,3,4\}$: $h_3=\min(3,2,4)=2$, $h_4=\min(2,3,1)=1$.
    \item $S_4$ has $1$'s at rows $\{0,2,3\}$: $h_3=\min(1,0,2)=0$, $h_4=\min(4,0,3)=0$.
\end{itemize}

The full signature matrix is
\[
\begin{array}{c|cccc}
 & S_1 & S_2 & S_3 & S_4\\\hline
h_1 & 1 & 3 & 0 & 1\\
h_2 & 0 & 2 & 0 & 0\\
h_3 & 1 & 0 & 2 & 0\\
h_4 & 3 & 0 & 1 & 0
\end{array}
\]
\end{proof}

%================================================================
\begin{homework}
 For the matrix in Figure 3 with $h_1(x)=2x+1\bmod 6,\;h_2(x)=3x+2\bmod 6,\;h_3(x)=5x+2\bmod 6$:
\begin{enumerate}
    \item[(a)] Compute the minhash signatures.
    \item[(b)] Which hashes are true permutations?
    \item[(c)] Compare estimated and true Jaccard similarities for the six pairs.
\end{enumerate}
\end{homework}

\begin{proof}[\solutionname]

\textbf{(a) Signatures.} The hash values are
\[
\begin{array}{c|cccccc}
x & 0 & 1 & 2 & 3 & 4 & 5\\\hline
h_1=2x+1\bmod 6 & 1 & 3 & 5 & 1 & 3 & 5\\
h_2=3x+2\bmod 6 & 2 & 5 & 2 & 5 & 2 & 5\\
h_3=5x+2\bmod 6 & 2 & 1 & 0 & 5 & 4 & 3
\end{array}
\]
The columns have $1$'s at rows
$S_1{=}\{2\}$, $S_2{=}\{0,1\}$, $S_3{=}\{2,3,4,5\}$, $S_4{=}\{0,2,5\}$. Taking $\min$ over those rows:
\[
\begin{array}{c|cccc}
 & S_1 & S_2 & S_3 & S_4\\\hline
h_1 & 5 & 1 & 1 & 1\\
h_2 & 2 & 2 & 2 & 2\\
h_3 & 0 & 1 & 0 & 0
\end{array}
\]

\textbf{(b) Permutations.}
$h_1$ takes values $\{1,3,5,1,3,5\}$, hitting only odd residues; not a permutation. $h_2$ takes values $\{2,5,2,5,2,5\}$; not a permutation. $h_3$ takes values $\{2,1,0,5,4,3\}$, all six residues exactly once; $h_3$ \emph{is} a true permutation. (We know $ax+b\bmod n$ is a permutation iff $\gcd(a,n)=1$, and $\gcd(5,6)=1$ while $\gcd(2,6)=\gcd(3,6)>1$.)

\textbf{(c) Estimated vs.\ true Jaccards.}
True Jaccards from the columns:
\[
\begin{aligned}
J(S_1,S_2)&=0/3=0, & J(S_1,S_3)&=1/4=0.25, & J(S_1,S_4)&=1/3,\\
J(S_2,S_3)&=0/6=0, & J(S_2,S_4)&=1/4=0.25, & J(S_3,S_4)&=2/5=0.4.
\end{aligned}
\]
Estimated Jaccards from the signature (fraction of hash rows on which the two columns agree):
\[
\begin{aligned}
\hat J(S_1,S_2)&=1/3\;(h_2 \text{ only}),\\
\hat J(S_1,S_3)&=2/3;(h_2,h_3),\\
\hat J(S_1,S_4)&=2/3\;(h_2,h_3),\\
\hat J(S_2,S_3)&=2/3\;(h_1,h_2),\\
\hat J(S_2,S_4)&=2/3\;(h_1,h_2),\\
\hat J(S_3,S_4)&=1.
\end{aligned}
\]

\emph{Possible Reason:} The estimates are systematically too high. Two reasons we can tell: (i) only $h_3$ is a true permutation; $h_1,h_2$ collapse rows into very few buckets ($\{1,3,5\}$ and $\{2,5\}$ respectively), inflating the chance of accidental agreement and thus over-estimating similarity; (ii) the sample size of $3$ hash functions is far too small to yield concentration around the true Jaccard. Using many independent true permutations would drive $\hat J\to J$ by the law of large numbers (each indicator $\bm 1\{h(S_i)=h(S_j)\}$ is Bernoulli with mean exactly $J(S_i,S_j)$).
\end{proof}

\end{document}
